\documentclass[blue]{beamer}
\usetheme{Madrid}
%\setbeamertemplate{footline}{Chris Walters (MIT)}
\usenavigationsymbolstemplate{}
\setbeamertemplate{navigation symbols}{}
\setbeamertemplate{footline}{}
\usepackage{graphicx}
\usepackage{listings}
\usepackage{graphics}
\usepackage{multirow}
\usepackage{xkeyval}
\usepackage{verbatim}
\usepackage{amsmath, amsthm, amssymb}
\usepackage{pdfpages}
\usepackage{xcolor}
\newsavebox\MBox
\newcommand\Cline[2][red]{{\sbox\MBox{$#2$}%
  \rlap{\usebox\MBox}\color{#1}\rule[-1.2\dp\MBox]{\wd\MBox}{0.75pt}}}
\newcommand{\Perp}{\perp \! \! \! \perp}


\makeatletter
\makeatother

\title{Self Selection}
\author[]{Raffaele Saggio}
\institute{\small{Econ 560}}

\begin{document}

{
\thispagestyle{empty}
\setbeamertemplate{footline}{}
\frame{\titlepage}
}
\frame{

\frametitle{Self-selection}
\begin{itemize}
\item {\footnotesize Classic idea in labor economics: }{\footnotesize\emph{Self-selection}}{\footnotesize\medskip{}
}{\footnotesize\par}
\item {\footnotesize Individuals choose between opportunities based on heterogeneous
unobserved returns \medskip{}
}{\footnotesize\par}
\item {\footnotesize We've already encountered some versions of this, e.g.
in Keane and Wolpin (1997) \medskip{}
}{\footnotesize\par}
\item {\footnotesize Applications: occupational choice, immigration, labor
force participation, educational choice, treatment choice, any economics
subfield... \medskip{}
}{\footnotesize\par}
\item {\footnotesize We'll start with general discussion of self-selection
models and some related econometrics, then look at some applications}{\footnotesize\par}
\end{itemize}
}

\frame{

\frametitle{Roy Model}
\begin{itemize}
\item {\footnotesize Roy (1951) sought to understand the influence of occupational
choice on the observed distribution of earnings}{\scriptsize\medskip{}
}{\scriptsize\par}
\item {\footnotesize Consider individuals indexed by $i$ choosing a binary
variable $D_{i}\in\{0,1\}$, e.g. hunting vs. fishing}{\scriptsize\medskip{}
}{\scriptsize\par}
\item {\footnotesize$Y_{i1}$ and $Y_{i0}$ are $i$'s potential outcomes
(e.g. earnings) with and without treatment}{\scriptsize\medskip{}
}{\scriptsize\par}
\item {\footnotesize Realized outcome is $Y_{i}=Y_{i0}+(Y_{i1}-Y_{i0})D_{i}$}{\scriptsize\medskip{}
}{\scriptsize\par}
\item {\footnotesize Pure Roy (1951) model: Individuals want to maximize
$Y_{i}$, so choose the alternative with the best potential outcome:}{\footnotesize\par}
\end{itemize}
\begin{center}
{\footnotesize$D_{i}=1\left\{ Y_{i1}>Y_{i0}\right\} $}{\footnotesize\par}
\par\end{center}

}

\frame{

\frametitle{Roy Model}
\begin{itemize}
\item {\footnotesize Questions of interest:\medskip{}
}{\footnotesize\par}
\begin{itemize}
\item {\footnotesize Will the best hunters hunt?\medskip{}
}{\footnotesize\par}
\item {\footnotesize Will the best fishermen/women fish?\medskip{}
}{\footnotesize\par}
\end{itemize}
\item {\footnotesize Suppose potential outcomes are given by}{\footnotesize\par}
\end{itemize}
\begin{center}
{\footnotesize$Y_{id}=p_{d}S_{id},\ d\in\{0,1\}$}{\footnotesize\par}
\par\end{center}
\begin{itemize}
\item {\footnotesize$S_{id}$ is skill in occupation $d$, and $p_{d}$
is price of output\medskip{}
}{\footnotesize\par}
\item {\footnotesize A worker who is indifferent between the two occupations
satisfies}{\footnotesize\par}
\end{itemize}
\begin{center}
{\footnotesize$\log S_{i1}=\log p_{0}-\log p_{1}+\log S_{i0}$}{\footnotesize\par}
\par\end{center}

}

{ \setbeamercolor{background canvas}{bg=} \includepdf{aux/roy_1}
}

\frame{

\frametitle{Roy Model: Special Cases}
\begin{itemize}
\item {\footnotesize Suppose there is no variation in the non-treated outcome,
so $S_{i0}=\bar{S}_{0}\ \forall i$\medskip{}
}{\footnotesize\par}
\item {\footnotesize In this case the decision rule is}{\footnotesize\par}
\end{itemize}
\begin{center}
{\footnotesize$D_{i}=1\left\{ S_{i1}\geq\left(\tfrac{p_{0}}{p_{1}}\right)\bar{S}_{0}\right\} $}{\footnotesize\par}
\par\end{center}
\begin{itemize}
\item {\footnotesize Those with the most skill in sector 1 choose $D_{i}=1$\medskip{}
}{\footnotesize\par}
\item {\footnotesize Everyone with $D_{i}=1$ earns more than anyone with
$D_{i}=0$}{\footnotesize\par}
\end{itemize}
}

{ \setbeamercolor{background canvas}{bg=} \includepdf{aux/roy_2}
}

\frame{

\frametitle{Roy Model: Special Cases}
\begin{itemize}
\item {\footnotesize Suppose we have perfect correlation between $\log S_{i0}$
and $\log S_{i1}$:}{\footnotesize\par}
\end{itemize}
\begin{center}
{\footnotesize$\log S_{i1}=\alpha_{0}+\alpha_{1}\log S_{i0},\ \alpha_{1}>0$}{\footnotesize\par}
\par\end{center}
\begin{itemize}
\item {\footnotesize This is a one-factor model\medskip{}
}{\footnotesize\par}
\item {\footnotesize Decision rule:}{\footnotesize\par}
\end{itemize}
\begin{center}
{\footnotesize$D_{i}=1\left\{ \alpha_{0}+\alpha_{1}\log S_{i0}\geq\log p_{0}-\log p_{1}+\log S_{i0}\right\} $}{\footnotesize\par}
\par\end{center}

\begin{center}
{\footnotesize$=1\left\{ (\alpha_{1}-1)\log S_{i0}\geq\log p_{0}-\log p_{1}-\alpha_{0}\right\} $}{\footnotesize\par}
\par\end{center}
\begin{itemize}
\item {\footnotesize Higher skilled choose $D_{i}=1$ iff $\alpha_{1}\geq1$\medskip{}
}{\footnotesize\par}
\item {\footnotesize Note that $Var(\log S_{i1})=\alpha_{1}^{2}Var(\log S_{i0})$
. Higher skilled choose the sector with higher variance of earnings}{\footnotesize\par}
\end{itemize}
}

{ \setbeamercolor{background canvas}{bg=} \includepdf{aux/roy_3}
}

{ \setbeamercolor{background canvas}{bg=} \includepdf{aux/roy_4}
}


\frame{

\frametitle{Generalized Roy Model}
\begin{itemize}
\item {\footnotesize Generalized Roy model (Eisenhauer et al., 2015): Preference
for alternative $d$ depends on $Y_{id}$ as well as a cost $C_{id}$:}{\footnotesize\par}
\end{itemize}
\begin{center}
{\footnotesize$D_{i}=1\left\{ Y_{i1}-C_{i1}>Y_{i0}-C_{i0}\right\} $}{\footnotesize\par}
\par\end{center}
\begin{itemize}
\item {\footnotesize Allows us to ask richer questions about selection on
both levels and gains\medskip{}
}{\footnotesize\par}
\begin{itemize}
\item {\footnotesize Is average $Y_{i1}$ higher or lower for individuals
that choose $D_{i}=1$?\medskip{}
}{\footnotesize\par}
\item {\footnotesize Is average $Y_{i0}$ higher or lower for individuals
that choose $D_{i}=1$?\medskip{}
}{\footnotesize\par}
\item {\footnotesize Are average }{\footnotesize\emph{gains}}{\footnotesize{}
$Y_{i1}-Y_{i0}$ larger or smaller for individuals that choose $D_{i}=1$?\medskip{}
}{\footnotesize\par}
\end{itemize}
\item {\footnotesize Close link between generalized Roy model and econometric
models of treatment effect heterogeneity}{\footnotesize\par}
\end{itemize}
}

\frame{

\frametitle{Generalized Roy Model}
\begin{itemize}
\item {\footnotesize Write outcomes and costs as functions of characteristics
$X_{i}$:}{\footnotesize\par}
\end{itemize}
\begin{center}
{\footnotesize$Y_{id}=\mu_{d}(X_{i})+\epsilon_{id},\ C_{id}=\varphi_{d}(X_{i})+\eta_{id},\ d\in\{0,1\}$}{\footnotesize\par}
\par\end{center}
\begin{itemize}
\item {\footnotesize$E[\epsilon_{id}|X_{i}]=E[\eta_{id}|X_{i}]=0$ by definition
of $\mu_{d}(\cdot)$ and $\varphi_{d}(\cdot)$\medskip{}
}{\footnotesize\par}
\item {\footnotesize Decision rule is}{\footnotesize\par}
\end{itemize}
\begin{center}
{\footnotesize$D_{i}=1\left\{ \mu_{1}(X_{i})+\epsilon_{i1}-\varphi_{1}(X_{i})-\eta_{i1}>\mu_{0}(X_{i})+\epsilon_{i0}-\varphi_{0}(X_{i})-\eta_{i0}\right\} $}{\footnotesize\par}
\par\end{center}

\begin{center}
{\footnotesize$\implies D_{i}=1\left\{ \psi(X_{i})>v_{i}\right\} $}{\footnotesize\par}
\par\end{center}
\begin{itemize}
\item {\footnotesize Here $\psi(X_{i})=[\mu_{1}(X_{i})-\mu_{0}(X_{i})]-[\varphi_{1}(X_{i})-\varphi_{0}(X_{i})]$,
and $v_{i}=[\eta_{i1}-\eta_{i0}]-[\epsilon_{i1}-\epsilon_{i0}]$\medskip{}
}{\footnotesize\par}
\item {\footnotesize Dependence of decision on potential outcomes is embedded
in $\psi(\cdot)$ and $v_{i}$}{\footnotesize\par}
\end{itemize}
}

\frame{

\frametitle{Roy Model Identification}
\begin{itemize}
\item {\scriptsize Consider the regression of $Y_{i}$ on $X_{i}$ in the
sample of individuals who chose $D_{i}=1$\medskip{}
}{\scriptsize\par}
\item {\scriptsize The conditional expectation function is\medskip{}
 
\[
E[Y_{i1}|X_{i},D_{i}=1]=\mu_{1}(X_{i})+E[\epsilon_{i1}|X_{i},D_{i}=1]
\]
\medskip{}
\[
=\mu_{1}(X_{i})+E[\epsilon_{i1}|X_{i},v_{i}<\psi(X_{i})]
\]
}{\scriptsize\par}
\item {\scriptsize Since $v_{i}$ includes $\epsilon_{i1}$, generally expect
$E[\epsilon_{i1}|X_{i},v_{i}<\psi(X_{i})]\neq0$, so observed mean
outcome does not reflect population mean potential outcome $\mu_{1}(X_{i})$\medskip{}
}{\scriptsize\par}
\item {\scriptsize Heckman (1974, 1976, 1979) idea: Work out the bias term
to correct for self-selection}{\scriptsize\par}
\end{itemize}
}

\frame{

\frametitle{Correcting Selection}
\begin{itemize}
\item {\footnotesize Suppose 
\[
(\epsilon_{id},v_{i})^{\prime}|X_{i}\sim N\left(0,\left[\begin{smallmatrix}\sigma^{2} & \rho_{d}\sigma_{d}\\
\rho_{d}\sigma_{d} & 1
\end{smallmatrix}\right]\right),\ d\in\{0,1\}
\]
}{\footnotesize\par}
\item {\footnotesize This implies
\[
\epsilon_{id}=\rho_{d}\sigma_{d}v_{i}+e_{id},\mbox{ }E[e_{id}|v_{i},X_{i}]=0
\]
}{\footnotesize\par}
\item {\footnotesize Conditional expectation function in selected sample
is 
\[
E[Y_{i1}|X_{i},D_{i}=1]=\mu_{1}(X_{i})+\rho_{1}\sigma_{1}E[v_{i}|X_{i},v_{i}<\psi(X_{i})]
\]
\[
=\mu_{1}(X_{i})+\rho_{1}\sigma_{1}\lambda_{1}\left(\psi(X_{i})\right)
\]
}{\footnotesize\par}
\item {\footnotesize$\lambda_{1}(c)\equiv-\frac{\phi(c)}{\Phi(c)}$ is the
inverse Mills ratio\medskip{}
 }{\footnotesize\par}
\item {\footnotesize Likewise $E[Y_{i0}|X_{i},D_{i}=0]=\mu_{0}(X_{i})+\rho_{0}\sigma_{0}\lambda_{0}(\psi(X_{i}))$,
with $\lambda_{0}(c)=\phi(c)/[1-\Phi(c)]$\medskip{}
}{\footnotesize\par}
\item {\footnotesize What restrictions on $(\rho_{1}\sigma_{1},\rho_{0}\sigma_{0})$
are implied by a pure Roy model?}{\footnotesize\par}
\end{itemize}
}

\frame{

\frametitle{Generalized Roy Model Identification}
\begin{itemize}
\item {\scriptsize Two-step ``Heckit'' approach: Estimate $\psi(X_{i})$
in a first-step probit, then include $\lambda_{d}(\hat{\psi}(X_{i}))$
in second-step OLS regression\medskip{}
}{\scriptsize\par}
\item {\scriptsize Mills ratio is a }{\scriptsize\textbf{control function}}{\scriptsize{}
that adjusts for selection bias \medskip{}
}{\scriptsize\par}
\item {\scriptsize Problem: Dependent on functional form \medskip{}
}

{\scriptsize\par}\begin{itemize}
\item {\scriptsize If $X_{i}$ is a constant, cannot estimate \medskip{}
}{\scriptsize\par}
\item {\scriptsize If $\mu(X_{i})$ and $\psi(X_{i})$ are unrestricted,
cannot estimate\medskip{}
}{\scriptsize\par}
\item {\scriptsize Estimable if we restrict $\mu_{d}(X_{i})$ and $\psi(X_{i})$
(e.g. make them linear) so that $\mu(X_{i})$ is not colinear with
$\lambda_{d}(\psi(X_{i}))$, but typically no good justification for
this}{\scriptsize\par}
\end{itemize}
\end{itemize}
}

\frame{

\frametitle{Generalized Roy Model Identification}
\begin{itemize}
\item {\footnotesize Solution: Suppose we have another variable $Z_{i}$
that affects costs but not potential outcomes:}{\footnotesize\par}
\end{itemize}
\begin{center}
{\footnotesize$E[Y_{id}|X_{i},Z_{i}]=\mu_{d}(X_{i}),\ E[C_{id}|X_{i},Z_{i}]=\varphi_{d}(X_{i},Z_{i})$}{\footnotesize\par}
\par\end{center}
\begin{itemize}
\item {\footnotesize Participation rule becomes }{\footnotesize\par}
\end{itemize}
\begin{center}
{\footnotesize$D_{i}=1\left\{ \psi(X_{i},Z_{i})>v_{i}\right\} $.}{\footnotesize\par}
\par\end{center}
\begin{itemize}
\item {\footnotesize Under normality, control function terms equal $\rho_{d}\sigma_{d}\lambda_{d}(\psi(X_{i},Z_{i}))$,
which vary conditional on $X_{i}$\medskip{}
}{\footnotesize\par}
\item {\footnotesize Sound familiar? }{\footnotesize\par}
\medskip
\item {\footnotesize More on this in Econ 628! }
\end{itemize}
}

\section{Mulligan/Rubinstein}

\frame{

\frametitle{Mulligan and Rubinstein (2008)}
\begin{itemize}
\item {\small Mulligan and Rubinstein (QJE 2008) apply the logic of self-selection
models to study female wages and labor force participation \medskip{}
}{\small\par}
\item {\small Since the 1970s, female labor force participation has grown
dramatically (from 46\% in 1975 to 59\% in 2007) \medskip{}
}{\small\par}
\item {\small The male/female wage gap has fallen, but within-gender inequality
has increased for females\medskip{}
}{\small\par}
\item {\small MR hope to make sense of this by studying changes in the pattern
of selection into the labor market \medskip{}
}{\small\par}
\item {\small Idea: Growing returns to human capital may produce more positive
selection bias, closing the gender wage gap but increasing within-gender
differences }{\small\par}
\end{itemize}
}

{ \setbeamercolor{background canvas}{bg=} \includepdf{aux/MR_f1}
}

\frame{

\frametitle{Mulligan and Rubinstein (2008)}
\begin{itemize}
\item MR begin with the potential log wage equation 
\[
w_{it}=\mu_{t}^{w}+g_{i}\gamma_{t}+\sigma_{t}^{w}\epsilon_{it}^{w}
\]
\item $\mu_{t}^{w}$ is the overall mean potential log wage at time $t$ 
\item $\sigma_{t}^{w}$ is the residual standard deviation of wages at time
$t$ 
\item $\gamma_{t}$ is the mean gender \emph{potential} wage gap 
\item \emph{Observed} wage gap $G_{t}$ (assume all men work): 
\[
G_{t}=E[w_{it}|g_{i}=1,L_{it}=1]-E[w_{it}|g_{i}=0]
\]
\[
=\gamma_{t}+\sigma_{t}^{w}E[\epsilon_{it}|g_{i}=1,L_{it}=1]
\]
\[
=\gamma_{t}+\sigma_{t}^{w}b_{t}.
\]
\end{itemize}
}

\frame{

\frametitle{Mulligan and Rubinstein (2008)}
\begin{itemize}
\item Change in measured wage gaps: 
\[
\Delta G_{t}=\Delta\gamma_{t}+b_{t-1}\Delta\sigma_{t}^{w}+\sigma_{t}^{w}\Delta b_{t}
\]
\item Components due to change in potential wage gap, change in inequality
for given selection bias, change in selection bias {\small\medskip{}
}{\small\par}
\item MR want to estimate the third term 
\end{itemize}
}

\frame{

\frametitle{Mulligan and Rubinstein (2008)}
\begin{itemize}
\item {\small Participation determined by reservation wage $r_{it}$: 
\[
r_{it}=\mu_{t}^{r}+\sigma_{t}^{r}\epsilon_{it}^{r}
\]
}{\small\par}
\item {\small Work decision is $L_{it}=1\left\{ r_{it}<w_{it}\right\} $,
so }{\small\par}
\end{itemize}
{\small
\[
L_{it}=1\{\epsilon_{it}^{r}-\tfrac{\sigma_{t}^{w}}{\sigma_{t}^{r}}\epsilon_{it}^{w}<\tfrac{\gamma_{t}+\mu_{t}^{w}-\mu_{t}^{r}}{\sigma_{t}^{r}}\}
\]
}{\small\par}
\begin{itemize}
\item {\small Behavioral change driven by factors common to all women (right-hand
side), or factors that interact with womens' characteristics (left-hand
side) \medskip{}
}{\small\par}
\item {\small With small enough $\sigma_{t}^{w}$ and $\Cov(\epsilon_{it}^{r},\epsilon_{it}^{w})<0$, working
women will have lower potential wages (selection driven by low reservation
wages) \medskip{}
}{\small\par}
\item {\small As $\sigma_{t}^{w}$ increases, participation increasingly
driven by high wage offers $\implies$ bias more positive }{\small\par}
\end{itemize}
}

{ \setbeamercolor{background canvas}{bg=} \includepdf{aux/MR_f2}
}

\frame{

\frametitle{Mulligan and Rubinstein (2008)}
\begin{itemize}
\item Under normality selection bias can be written 
\[
b_{t}=\theta\left(\sigma_{t}^{w}/\sigma_{t}^{r}\right)\lambda\left(P_{t}\right)
\]
\item $P_{t}$ is the fraction of women who work; $\sigma_{t}^{w}/\sigma_{t}^{r}$
is the relative importance of market wages in the participation decision
{\small\medskip{}
}{\small\par}
\item Changes in bias can come from either source {\small\medskip{}
}{\small\par}
\item For a fixed selection rule, increasing $P_{t}$ reduces bias; but
an increase in $\sigma_{t}^{w}$ changes the selection rule, possibly
increasing bias 
\end{itemize}
}

\frame{

\frametitle{Mulligan and Rubinstein (2008)}
\begin{itemize}
\item MR assess changes in selection using the March CPS {\small\medskip{}
}{\small\par}
\item Two-step approach: 
\[
w_{it}=X_{it}\beta_{t}+g_{i}\gamma_{t}+g_{i}\theta_{t}\lambda(Z_{it}\delta_{t})+u_{it}
\]
\item $X$: Education, marital status, potential experience, region {\small\medskip{}
}{\small\par}
\item $Z$: Same as $X$, plus number of young children...
\end{itemize}
}

{ \setbeamercolor{background canvas}{bg=} \includepdf{aux/MR_t1}
}

{ \setbeamercolor{background canvas}{bg=} \includepdf{aux/MR_f3}
}

\frame{

\frametitle{Mulligan and Rubinstein (2008)}
\begin{itemize}
\item The OLS gender wage gap fell by 16 log points from 1975-1995 {\small\medskip{}
}{\small\par}
\item Controlling for the Mills ratio, the gap \emph{widened} {\small\medskip{}
}{\small\par}
\item Two-step estimates suggest a reversal of bias, from negative in the
1970s to positive in the 1990s {\small\medskip{}
}{\small\par}
\item Taken at face value, the results imply that all of the decrease in
the gender wage gap was due to changing selection 
\end{itemize}
}

{ \setbeamercolor{background canvas}{bg=}

\section{Chandra and Staiger}

\frame{

\frametitle{Chandra and Staiger (2007)}
\begin{itemize}
\item {\small Chandra and Staiger (JPE 2007) consider self selection in a
different context: treatment choice in health care \medskip{}
}{\small\par}
\item {\small Stylized fact: Large differences in medical treatments/spending
across areas, without commensurate differences in outcomes (see Dartmouth
map, or Atul Gawande) \medskip{}
}{\small\par}
\item {\small\textbf{Flat of the curve }}{\small explanation: High-use areas
perform a lot of expensive treatments with low returns, so big gains
could be had by cutting back \medskip{}
}{\small\par}
\item {\small CS ask whether this interpretation is compromised by self-selection }{\small\par}
\end{itemize}
}

{ \setbeamercolor{background canvas}{bg=} \includepdf{aux/skinner_1}
}

\frame{

\frametitle{Chandra and Staiger (2007)}
\begin{itemize}
\item {\footnotesize CH model: Treatment choice with productivity spillovers
\medskip{}
}{\footnotesize\par}
\item {\footnotesize Consider the case of heart attacks (acute myocardial
infarction, AMI) \medskip{}
}{\footnotesize\par}
\item {\footnotesize Two possible treatments: \medskip{}
}

\begin{enumerate}
\item {\footnotesize Non-intensive medical management \medskip{}
}{\footnotesize\par}
\item {\footnotesize Intensive intervetion (surgery; bypass or angioplasty)
\medskip{}
}{\footnotesize\par}
\end{enumerate}
\item {\footnotesize Physician chooses treatment option $i\in\{1,2\}$ that
maximizes utility for each patient, which depends on survival probability
$Survival_{i}$ and cost $Cost_{i}$ \medskip{}
}{\footnotesize\par}
\item {\footnotesize Survival probability depends positively on population
share receiving the same treatment, $P_{i}$ $\implies$ positive
productivity externality \medskip{}
}{\footnotesize\par}
\item {\footnotesize Justifications: Learning by doing (``see one, do one,
teach one''), availability of support services for commonly used
treatments, physician sorting }{\footnotesize\par}
\end{itemize}
}

\frame{

\frametitle{Chandra and Staiger (2007)}
\begin{itemize}
\item Formal Roy model: For $i\in\{1,2\}$ 
\[
Survival_{i}=\beta_{i}^{s}Z+\alpha_{i}^{s}P_{i}+\epsilon_{i}^{s}
\]
\[
Cost_{i}=\beta_{i}^{c}Z+\alpha_{i}^{c}P_{i}+\epsilon_{i}^{c}
\]
\[
U_{i}=Survival_{i}-\lambda Cost_{i}
\]
\[
=\beta_{i}Z+\alpha_{i}P_{i}+\epsilon_{i}
\]
\item Intensive treatment probability 
\[
Pr[i=2]=Pr[U_{2}>U_{1}]
\]
\[
=Pr[(\beta_{2}-\beta_{1})Z+\alpha_{2}P_{2}-\alpha_{1}(1-P_{2})+(\epsilon_{2}-\epsilon_{1})>0]
\]
\[
=Pr[\alpha P_{2}-\alpha_{1}+\beta Z>\epsilon]
\]
\end{itemize}
}

\frame{

\frametitle{Chandra and Staiger (2007)}
\begin{itemize}
\item Equilibrium: $P_{2}$ must equal the fraction of individuals for whom
treatment 2 is chosen (fixed point) 
\[
P_{2}=\int Pr[\alpha P_{2}-\alpha_{1}+\beta Z>\epsilon]dF(Z)
\]
\[
=G(P_{2})
\]
\item Equilibrium proportion of intensive treatments must generate external
benefits that justify treating this proportion intensively {\small\medskip{}
}{\small\par}
\item Produces the possibility of multiple equilibria 
\end{itemize}
}

{ \setbeamercolor{background canvas}{bg=} \includepdf{aux/CS_f1b}
}

{ \setbeamercolor{background canvas}{bg=} \includepdf{aux/CS_f1a}
}

{ \setbeamercolor{background canvas}{bg=} \includepdf{aux/CS_f2a}
}

{ \setbeamercolor{background canvas}{bg=} \includepdf{aux/CS_f2b}
}

\frame{

\frametitle{Chandra and Staiger (2007)}
\begin{itemize}
\item {\scriptsize Implications of the CS model: \medskip{}
}

{\scriptsize\par}\begin{itemize}
\item {\scriptsize Possibly multiple equilibia \medskip{}
}{\scriptsize\par}
\item {\scriptsize Large differences in treatment rates can be generated
by small differences in patient characteristics \medskip{}
}{\scriptsize\par}
\item {\scriptsize Identical patients may be treated differently depending
on population treatment rate \medskip{}
}{\scriptsize\par}
\item {\scriptsize Gains to intensive treatment are larger for individuals
who are most likely to be treated intensively ($TOT>ATE$) \medskip{}
}{\scriptsize\par}
\item {\scriptsize Being in an intensive area is good for those most appropriate
for the intensive treatment, and worse for those least appropriate
\medskip{}
}{\scriptsize\par}
\item {\scriptsize Marginal patient receiving the intensive treatment in
an intensive area will be less appropriate for this treatment \medskip{}
}{\scriptsize\par}
\item {\scriptsize For a given patient, treatment effect of intensive treatment
is larger in an intensive area \medskip{}
}{\scriptsize\par}
\end{itemize}
\item {\scriptsize This model can explain why we might see large differences
in treatment intensity without large differences in average outcomes }{\scriptsize\par}
\end{itemize}
}

\frame{

\frametitle{Chandra and Staiger (2007)}
\begin{itemize}
\item {\small Welfare (expected utility): If differences in intensity are
generated by multiple equilibria w/identical patients, can rank equilibria
by computing costs and outcomes \medskip{}
}{\small\par}
\item {\small In this case, social planner will want to reduce variation
in intensity (switch to less intense equilibrium) if intense areas
have higher costs and no better outcomes \medskip{}
}{\small\par}
\item {\small If there is a single equilibrium with intensity that varies
because of differences in patient characteristics, }{\small\emph{not
enough}}{\small{} variation in intensity -- social planner would be
willing to harm marginal patient to increase utility of inframarginal
patients, and by assumption, decisionmaker (physician/patient) is
not. }{\small\par}
\end{itemize}
}

\frame{

\frametitle{Chandra and Staiger (2007)}
\begin{itemize}
\item {\footnotesize CS try to test the implications of their model \medskip{}
}{\footnotesize\par}
\item {\footnotesize Focus on heart attacks. Appropriate because: \medskip{}
}

\begin{itemize}
\item {\footnotesize A leading cause of death in the US \medskip{}
}{\footnotesize\par}
\item {\footnotesize Geographically distinct markets \medskip{}
}{\footnotesize\par}
\item {\footnotesize Two starkly different treatments: Intensive surgery
(bypass), non-intensive management (thrombolytics) \medskip{}
}{\footnotesize\par}
\item {\footnotesize Large differences across markets in intensity, but also
variation in appropriateness across patients \medskip{}
}{\footnotesize\par}
\end{itemize}
\item {\footnotesize CS look at the effects of intensive treatment using
data from the Cooperative Cardiovascular Project (CCP), a data set
collected to try to analyze quality of heart attack treatment }{\footnotesize\par}
\end{itemize}
}

\frame{

\frametitle{Chandra and Staiger (2007)}
\begin{itemize}
\item {\footnotesize Model for intensive treatment receipt: 
\[
Pr[\mbox{Cardiac Cath}_{ij}=1|X_{i},j]=G(\theta_{0}+\theta_{j}+X_{i}\Phi)
\]
}{\footnotesize\par}
\item {\footnotesize Now $i$ is patient, $j$ is hospital referral region
\medskip{}
}{\footnotesize\par}
\item {\footnotesize$X_{i}$ is detailed set of characteristics from patient
chart \medskip{}
}{\footnotesize\par}
\item {\footnotesize Fitted value $\hat{G}(\theta_{0}+X_{i}\Phi)$ captures
appropriateness after taking out HRR effect \medskip{}
}{\footnotesize\par}
\item {\footnotesize$\theta_{j}$ captures differences in treatment rates
for similar patients \medskip{}
}{\footnotesize\par}
\item {\footnotesize CS begin by testing for the presence of Roy selection:
Are gains to intensive treatment increasing in the odds of getting
this treatment? \medskip{}
}{\footnotesize\par}
\item {\footnotesize Model for effect of intensive treatment: 
\[
Y_{ij}=\beta_{0}+\beta_{1}\mbox{Cardiac Cath}_{ij}+X_{i}\Pi+u_{ijk}
\]
}{\footnotesize\par}
\item {\footnotesize Instrument using differential distance to nearest cath
hospital, splitting sample by $\hat{G}$ }{\footnotesize\par}
\end{itemize}
}

{ \setbeamercolor{background canvas}{bg=} \includepdf{aux/CS_t2}
}

{ \setbeamercolor{background canvas}{bg=} \includepdf{aux/CS_t1}
}

{ \setbeamercolor{background canvas}{bg=} \includepdf{aux/CS_f3}
}

\frame{

\frametitle{Chandra and Staiger (2007)}

\begin{small} 
\begin{itemize}
\item {\small Next, look for evidence of productivity spillovers \medskip{}
}

\begin{itemize}
\item {\small Is the quality of medical management worse in areas that do
less of it? \medskip{}
}{\small\par}
\item {\small Are identical individuals more likely to get intensive treatment
in areas where other patients are more appropriate for it? \medskip{}
}{\small\par}
\item {\small Is the return to intensive treatment higher in higher-intensity
areas? Roy model with spillovers says yes -- flat-of-the-curve explanation
says no \medskip{}
}{\small\par}
\item {\small Are the most-appropriate patients better off in high-intensity
regions? Are the least-appropriate patients worse off? }{\small\par}
\end{itemize}
\end{itemize}
\end{small} 

}

{ \setbeamercolor{background canvas}{bg=}\includepdf{aux/CS_4}
}

{ \setbeamercolor{background canvas}{bg=} \includepdf{aux/CS_t6}
}

{ \setbeamercolor{background canvas}{bg=} \includepdf{aux/CS_t7}
}

\frame{

\frametitle{Chandra and Staiger (2007)}
\begin{itemize}
\item {\footnotesize Many predictions of the CS model with spillovers have
support in the data \medskip{}
}{\footnotesize\par}
\item {\footnotesize Suggests that differences in treatment intensity across
areas may be due to optimal decision-making in the presence of productivity
spillovers, rather than inefficient utilization of care \medskip{}
}{\footnotesize\par}
\item {\footnotesize At the same time, welfare implications unclear \medskip{}
}

\begin{itemize}
\item {\footnotesize If patients are identical across areas, larger costs
with the same outcomes implies that expected utility could be increased
by reducing treatment intensity, regardless of whether variation is
caused by productivity spillovers/multiple equilibria or flat of the
curve \medskip{}
}{\footnotesize\par}
\item {\footnotesize If source of variation in intensity is patient heterogeneity,
we need }{\footnotesize\emph{more}}{\footnotesize{} specialization \medskip{}
}{\footnotesize\par}
\end{itemize}
\item {\footnotesize Not obvious how to test whether variation is due to
single vs. multiple equilibria }{\footnotesize\par}
\end{itemize}
}


\section{References}

\frame{ 

\frametitle{References}
\begin{itemize}
\item {\scriptsize Abramitzky, Ran, Leah Boustan, and Katherine Ericksson
(2012). ``Europe's Tired, Poor, Huddled Masses: Self-selection and
Economic Outcomes in the Age of Mass Migration.'' }{\scriptsize\emph{American
Economic Review }}{\scriptsize 102.}{\scriptsize\par}
\item {\scriptsize Bjorklund, Anders, and Robert Moffitt (1987). ``The Estimation
of Wage Gains and Welfare Gains in Self-Selection Models.'' }{\scriptsize\emph{Review
of Economics and Statistics }}{\scriptsize 69.}{\scriptsize\par}
\item {\scriptsize Borjas, George (1987). ``Self-selection and the Earnings
of Immigrants.'' }{\scriptsize\emph{American Economic Review }}{\scriptsize 77.}{\scriptsize\par}
\item {\scriptsize Brinch, Christian, Magne Mogstad, and Matthew Wiswall
(2017). ``Beyond LATE with a Discrete Instrument.'' }{\scriptsize\emph{Journal
of Political Economy }}{\scriptsize 125.}{\scriptsize\par}
\item {\scriptsize Carneiro, Pedro, James Heckman, and Edward Vytlacil (2010).
``Evaluating Marginal Policy Changes and the Average Effect of Treatment
for Individuals at the Margin.'' }{\scriptsize\emph{Econometrica
}}{\scriptsize 78.}{\scriptsize\par}
\item {\scriptsize Chandra, Amitabh and Douglas Staiger (2007). \textquotedblleft Productivity
Spillovers in Health Care: Evidence from the Treatment of Heart Attacks.\textquotedblright{}
}{\scriptsize\textit{Journal of Political Economy}}{\scriptsize{} 115.}{\scriptsize\par}
\item {\scriptsize Chiquiar, Daniel and Gordon Hanson (2005). ``International
Migration, Self-selection, and the Distribution of Wages: Evidence
from Mexico and the United States.'' }{\scriptsize\emph{Journal of
Political Economy }}{\scriptsize 113.}{\scriptsize\par}
\end{itemize}
}

\frame{ 

\frametitle{References}
\begin{itemize}
\item {\scriptsize Cornelissen, Thomas, Christian Dustmann, Anna Raute, and
Uta Schonberg (forthcoming). ``Who Benefits From Universal Child
Care? Estimating Marginal Returns to Early Child Care Attendance.''
}{\scriptsize\emph{Journal of Political Economy.}}{\scriptsize\par}
\item {\scriptsize Dahl, Gordon (2002). \textquotedblleft Mobility and the
Return to Education: Testing a Roy Model with Multiple Markets.\textquotedblright{}
}{\scriptsize\textit{Econometrica}}{\scriptsize{} 70.}{\scriptsize\par}
\item {\scriptsize Eisenhauer, Philipp, James Heckman, and Edward Vytlacil
(2015). ``The Generalized Roy Model and the Cost-Benefit Analysis
of Social Programs.'' }{\scriptsize\emph{Journal of Political Economy
}}{\scriptsize 123.}{\scriptsize\par}
\item {\scriptsize French, Eric, and Christopher Taber (2011). ``Identification
of Models of the Labor Market.'' }{\scriptsize\uline{Handbook of
Labor Economics}}{\scriptsize{} Volume 4A.}{\scriptsize\par}
\item {\scriptsize Gronau, Reuben (1974). ``Wage Comparisons - A Selectivity
Bias.'' }{\scriptsize\emph{Journal of Political Economy }}{\scriptsize 82.}{\scriptsize\par}
\item {\scriptsize Heckman, James (1974): \textquotedblleft Shadow prices,
market wages, and labor supply,\textquotedblright{} }{\scriptsize\emph{Econometrica}}{\scriptsize{}
42.}{\scriptsize\par}
\item {\scriptsize Heckman, James (1976). ``The Common Structure of Statistical
Models of Truncation, Sample Selection and Limited Dependent Variables
and a Simple Estimator for Such Models.'' }{\scriptsize\emph{Annals
of Economic and Social Measurement }}{\scriptsize 5.}{\scriptsize\par}
\end{itemize}
}

\frame{ 

\frametitle{References}
\begin{itemize}
\item {\scriptsize Heckman, James (1979). ``Sample Selection Bias as a Specification
Error.'' }{\scriptsize\emph{Econometrica }}{\scriptsize 47.}{\scriptsize\par}
\item {\scriptsize Heckman, James and Bo Honor� (1990). \textquotedblleft The
Empirical Content of the Roy Model.\textquotedblright{} }{\scriptsize\textit{Econometrica}}{\scriptsize{}
58.}{\scriptsize\par}
\item {\scriptsize Heckman, James, Sergio Urzua, and Edward Vytlacil (2006).
``Understanding Instrumental Variables in Models with Essential Heterogeneity.''
}{\scriptsize\emph{Review of Economics and Statistics 88.}}{\scriptsize\par}
\item {\scriptsize Heckman, James, and Edward Vytlacil (2005). ``Structural
Equations, Treatment Effects, and Econometric Policy Evaluation.''
}{\scriptsize\emph{Econometrica }}{\scriptsize 73.}{\scriptsize\par}
\item {\scriptsize Heckman, James, and Edward Vytlacil (2007). ``Econometric
Evaluation of Social Programs, Part II: Using the Marginal Treatment
Effect to Organize Alternative Economic Estimators to Evaluate Social
Programs and to Forecast Their Effects in New Environments.'' }{\scriptsize\uline{Handbook
of Econometrics}}{\scriptsize{} Volume 6B.}{\scriptsize\par}
\item {\scriptsize Imbens, Guido, and Joshua Angrist (1994). ``Identification
and Estimation of Local Average Treatment Effects.'' }{\scriptsize\emph{Econometrica
}}{\scriptsize 62.}{\scriptsize\par}
\item {\scriptsize Kirkeboen, Lars, Edwin Leuven, and Magne Mogstad (2016).
``Field of Study, Earnings, and Self-selection.'' }{\scriptsize\emph{Quarterly
Journal of Economics }}{\scriptsize 131.}{\scriptsize\par}
\end{itemize}
}

\frame{ 

\frametitle{References}
\begin{itemize}
\item {\scriptsize Kline, Patrick, and Christopher Walters (2016). ``Evaluating
Public Programs with Close Substitutes: The Case of Head Start.''
}{\scriptsize\emph{Quarterly Journal of Economics }}{\scriptsize 132.}{\scriptsize\par}
\item {\scriptsize Kline, Patrick, and Christopher Walters (2018). ``On
Heckits, LATE, and Numerical Equivalence.'' Working paper.}{\scriptsize\par}
\item {\scriptsize Mulligan, Casey and Yona Rubinstein (2008). \textquotedblleft Selection,
Investment, and Women\textquoteright s Relative Wages Over Time.\textquotedblright{}
}{\scriptsize\textit{Quarterly Journal of Economics}}{\scriptsize{}
123.}{\scriptsize\par}
\item {\scriptsize Olsen, Randall (1980). ``A Least Squares Correction for
Selectivity Bias.'' }{\scriptsize\emph{Econometrica }}{\scriptsize 48.}{\scriptsize\par}
\item {\scriptsize Rothschild, Casey and Florian Scheuer (2013). ``Redistributive
Taxation in the Roy Model.'' }{\scriptsize\emph{Quarterly Journal
of Economics }}{\scriptsize 123.}{\scriptsize\par}
\item {\scriptsize Roy, Andrew (1951). \textquotedblleft Some Thoughts on
the Distribution of Earnings.\textquotedblright{} }{\scriptsize\textit{Oxford
Economic Papers New Series}}{\scriptsize{} 3.}{\scriptsize\par}
\item {\scriptsize Vytlacil, Edward (2002). ``Independence, Monotonicity,
and Latent Index Models: An Equivalence Result.'' }{\scriptsize\emph{Econometrica
}}{\scriptsize 70.}{\scriptsize\par}
\end{itemize}
}

\frame{ 

\frametitle{References}
\begin{itemize}
\item {\scriptsize Willis, Robert and Sherwin Rosen (1979). \textquotedblleft Education
and Self-selection,\textquotedblright{} }{\scriptsize\textit{Journal
of Political Economy }}{\scriptsize 87.}{\scriptsize\par}
\item {\scriptsize Walters, Christopher (forthcoming). ``The Demand for
Effective Charter Schools.'' }{\scriptsize\emph{Journal of Political
Economy}}{\scriptsize .}{\scriptsize\par}
\end{itemize}
}
\end{document}